This is the file capable of actually instrument the target binary.
It uses two libraries :
\begin{itemize}
    \item \textbf{lief} : python library used mainly to add the dynamic dependency, the new sections and the relocation entries. 
        This library is used for the Binary Modification phase
    \item \textbf{GRAPE} : binary rewriting framework used to instrument every \gls{IP} found.
        This framework is used in the Instrumentation Phase
\end{itemize}
\subsection{Binary Modification}
Before actually instrument the binary the previously mentioned modifications must be made.
The principal steps are summarized in \cref{fig:binary_modification}.
\begin{figure}[ht]
    \centering
    \includegraphics[width=400pt]{BinaryModification.png}
    \caption{Binary modification phase in which the needed sections, relocations and dynamic dependencies are added. (\small The dotted arrows specify which informations are needed for the relocation phase) }
    \label{fig:binary_modification}
\end{figure}\\

\begin{center}
\textbf{create .my\_got section}    
\end{center}
The created section has 
\begin{itemize}
    \item \textit{section type = PROGBITS} . This means that the section data are actually into the \gls{ELF} file.
        For example the \textit{.bss} section is declared as NOBITS because its data are not stored into the \gls{ELF} but a 
        memory region must be allocated when the program starts
    \item \textit{section flags} \begin{itemize}
        \item ALLOC : The section must occupy actual space in memory. This memory must be allocated during \gls{ELF} loading.
        \item WRITE : This section must be writable by the process.
    \end{itemize}
    \item \textit{section content} : This section must contains the resolved addresses of the two symbols that are relocated.
        So the content of this section is composed of 16 bytes (in 64 bit addressing each resolved symbol occupy 8 bytes).
        Automatically the section content is initializated to 0.
    \item \textit{section alignment = 8} . Since this section will store 8 bytes addresses, the alignment is setted to 8.
\end{itemize}
This section, as stated before, must be actually loaded into memory. But this means associate the section to a segment 
with the LOAD flag, that has the same permits of the section.
\textbf{lief}, if specified, automatically add a new segment for each added section.

\begin{center}
\textbf{create symbols}    
\end{center}
Two symbols must be created. These two symbols correspond to the two exported symbols of the \gls{SO} and so the name must 
match. Specifically these symbols are created with:
\begin{itemize}
    \item \textit{symbol name} : Since the symbols name must match, the two symbols name are \begin{itemize}
        \item \textit{\_\_afl\_area\_ptr}
        \item \textit{\_\_forkserver\_start\_ptr}
    \end{itemize}
    \item \textit{symbol type = OBJECT}. The symbol represent a variable and not a functioning
    \item \textit{symbol binding = GLOBAL}. The symbol could be referenced by the dynamic linker. \\As stated in 
        \cref{sec:loading_relocations}, during \gls{ELF} loading, the dynamic linker resolve all the symbols for which there 
        is a relocation entry. A symbol could be exported or imported only if its binding is configured to be GLOBAL.
        There are LOCAL symbols that, by definition, are local to the object that define it and cannot be exported 
        (and therefore not imported either). 
    \item \textit{symbol visibility = DEFAULT}. The symbol has the default visibility of \gls{ELF}. There are other types 
        (such as HIDDEN or PROTECTED) that are used for local symbols. 
    \item \textit{symbol value = 0}. The symbol does not have a value. This is a clue to the fact that this symbol 
    must be resolved by the dynamic linker.
    \item \textit{symbol size = 8}. The symbol size is 8 bytes.
\end{itemize} 

\begin{center}
\textbf{Add relocation entries}    
\end{center}
Two relocation entries must be added, The first is used for resolving \textit{\_\_afl\_area\_ptr} and the second for 
\textit{\_\_forkserver\_start\_ptr}.
In order to create a relocation entry must be specified:
\begin{itemize}
    \item \textit{target addres of relocation} : where the resolved value will be written
    \item \textit{relocation type} : for both the symbols is used X86\_64\_GLOB\_DAT (that is the usual type when dealing with \gls{GOT} resolution)
    \item \textit{relocation encoding} : for both the symbols is used RELA encoding, that consist in esplicitly specify the addend.
    \item \textit{relocation porpouse} : for both the symbols is used a DYNAMIC porpouse because the symbol must be resolved at runtime bu the dynamic linker.
    \item \textit{addend = 0}
    \item \textit{relocation symbol} :the symbol to be resolved
\end{itemize}
The difference between the two relocation entries lies in the target address and in the relocation symbol.
For \textit{\_\_afl\_area\_ptr} the target address is the virtual address of \textit{.my\_got} section.\\
\textit{\_\_forkserver\_start\_ptr} the target address is the virtual address of \textit{.my\_got} section + 8.

\begin{center}
\textbf{Create .mydata section}    
\end{center}
All the consideration made for the previous section are already valid, a part from the section content.
In fact AFL, to know if a target has been instrumented or not, before actually run it, check if in some part of the entire 
\gls{ELF} is contained the string ``\_\_AFL\_SHM\_ID''. If is not contained AFL stops.\footnote{
    This check is actually avoidable setting the environmental variable AFL\_SKIP\_BIN\_CHECK=1 as must be done using ZAFL.
}
For this reason the aforementioned string is placed inside this section.
The string will be overwritten by the other data, but at fuzzing setup will pass the AFL check.
According to mine implementation this section will contains different informations at different addresses:
\begin{itemize}
    \item \textit{.mydata+0x10} : contains the ID of the previous \gls{BB}
    \item \textit{.mydata+0x20} : contains the address of \gls{SHM} initialized by the initial instrumentation
    \item \textit{.mydata+0x30} : contains the saved \texttt{rax} register to maintain its value before and after the instrumentation
\end{itemize}
At the end of the \textbf{Binary Modification phase}, the modified binary is saved to the filesystem.
\subsection{Instrumentation}
The instrumentation is made entirely using GRAPE framework.
It takes care of adjusting all the references shift caused by the altered \textit{.text} section, alongside the eventual 
section size adjusting (caused by the exceeding in dimension of the section given the added instructions) and many other features.
GRAPE first reconstruct the \gls{CFG} of the program and then performs operations on the obtained \gls{CFG}.
The classes that the developed tool will use are:
\begin{itemize}
    \item \textsc{AddInstructionChange} : Used to add an instruction into a specific address of the binary \gls{CFG}.
        By default the instruction is added specific at the address, shifting the instruction that was there before, after.
    \item \textsc{RemoveInstructionChange} : Used to remove an instruction from a specific address of the binary
\end{itemize}
These two classes takes care of \gls{CFG} modification (split a \gls{BB}, union two \gls{BB}, \dots) but does not take care of RIP-relative instructions 
modification, togheter with symbol coherence.
An example of an error due to not taking care of modifying single instructions is \cref{fig:rip_relative_error_after_shift}.\\
\begin{figure}[ht]
    \centering
    \begin{subfigure}{0.8\textwidth}
        \centering
        \includegraphics[width=\textwidth]{error_rip_relative1.png}
        \caption{Example of a Basic Block before the addition of an instruction}
        \label{fig:rip_relative_error_after_shift1}
    \end{subfigure}
    \begin{subfigure}{0.8\textwidth}
        \centering
        \includegraphics[width=0.9\textwidth]{error_rip_relative2.png}
        \caption{Same Basic Block of \cref{fig:rip_relative_error_after_shift1} with a nop instruction added}
        \label{fig:rip_relative_error_after_shift2}
    \end{subfigure}
    \caption{Example of the error caused by not taking care of RIP-relative instruction offset}
    \label{fig:rip_relative_error_after_shift}
\end{figure}\footnote{
    As can be seen the jmp at 0x11ef seem to be shifted, but in reality the jmp instruction is always RIP-relative and
    the disassembler used \cite{cutter} perform automatically the relative address resolution.
}
In \cref{fig:rip_relative_error_after_shift} can be seen that, after having inserted a \texttt{nop} instruction, the \texttt{lea} 
that was at the address 0x11db now shift at the address 0x11dc.
For this reason the referenced location is no longer a valid data location. 
Before the addition the computation was $0x11db+0xe22+7 = 0x2004 <usage\_string\_location>$\footnote{
    The +7 is due to the instruction size
}.
Now the computation is $0x11dc+0xe22+7 = 0x2005 <??>$

For these reasons the last step of binary rewriting (after actually add all the instructions to the \gls{BB}) uses two other classes:
\begin{itemize}
    \item \textsc{RIPDataCoherenceAnalysis} : adjust the RIP-relative instructions.
    \item \textsc{SymbolCoherenceAnalysis} : adjus symbols. For example, after having insert some instruction, 
        the target address of another instruction has shifted. The reconstruction of the correct target address is
        performed by the SymbolCoherence phase
\end{itemize}
Since GRAPE takes care of all the shifting and adjusting, the instrumentation algorithm must only takes care 
of computing the right offset of the RIP-relative instructions to be instrumented and insert the instructions.

Since the instructions must be added to a specific address, all the offsets are computed as $final\_address - instrumentation\_address -instruction\_size$.
For example,with reference to \cref{lst:init_instr}, supposing the instrumentation address is 0x4100000c, the offset\_to\_previous\_saved\_ID can be computed as follows\\
$offset\_to\_previous\_saved\_ID=(\textit{.mydata}.virtual\_addr + 0x10) - 0x4100000c - 9$\footnote{
    The -9 at the end is needed because the instruction at line 6 of \cref{lst:init_instr} uses 9 bytes
}.

In \cref{alg:initial_instrumentation} can be seen the initial instrumentation algorithm.
As stated before, GRAPE automatically increment the size of \textit{.text} section if, adding instruction, 
the section is too small to contains them all.
When this happens the \textit{.mydata} and \textit{.my\_got} section could be shifted forward, and the previously computed 
offsets for the RIP-relative instructions are no longer valid. 
For this reason if a shift has been detected, recompute all the instructions.
Moreover can be seen that the instructions are inserted in reverse order (the index of the cycle starts from the last instruction).
This is because the instructions are inserted one on the top of the other.

\begin{algorithm}[ht]
    \caption{Initial Instrumentation algorithm}
    \label{alg:initial_instrumentation}
    \begin{algorithmic}[1]
        \REQUIRE $addr , mydata\_addr , mygot\_addr$ 
        \ENSURE all the initial instructions are inserted in the correct order at address ``addr'' 

        \STATE $instructions \gets \textsc{computeInitialInstrumentation}(addr,mydata\_addr,mygot\_addr)$ \COMMENT{\small \textcolor{green!60!black}{\textsc{computeInitialInstrumentation} compute the right offset and return a list of assembly instructions written as strings}}

        \FOR{$i=\textsc{len}(instructions)$ \TO $1$} 
                \STATE $instr \gets \textsc{asm}(instructions[\,i\,])$ \COMMENT{\small \textcolor{green!60!black}{\textsc{asm} function convert the assembly instruction passed as parameter as string into machine code}}
                \STATE $\textsc{AddInstructionChange}(addr,instr).\textsc{modify}(binary.cfg)$
                \IF{$\text{.mydata section shifted after .text enlargement}$}
                    \STATE $mydata\_addr \gets mydata.virtual\_address$
                    \STATE $mygot\_addr \gets mygot.virtual\_address$
                    \STATE $instructions \gets \textsc{computeInitialInstrumentation}(
                            addr,mydata\_addr,mygot\_addr)$ \COMMENT{\small \textcolor{green!60!black}{recompute all the instructions with the new offset. For the instructions already added GRAPE will takes care of the shift}}
                    \STATE $instr \gets \textsc{asm}(instructions[\,i\,])$
                    \STATE $\textsc{RemoveInstructionChange}(addr).\textsc{modify}(binary.cfg)$ \COMMENT{\small \textcolor{green!60!black}{remove and reinsert the last instruction in order to make sure that the offset to mydata section is correct}}
                    \STATE $\textsc{AddInstructionChange}(addr,instr).\textsc{modify}(binary.cfg)$
                \ENDIF
        \ENDFOR
    \end{algorithmic}
\end{algorithm}
\begin{center}
\textbf{Initial Instrumntation}    
\end{center}
\begin{lstlisting}[caption={Initial Instrumentation},label=lst:init_instr]
push rsi
push rdi
mov rax, QWORD PTR[rip + offset_to_relocated_forkserver_function]
mov rax, [rax]
call rax
mov WORD PTR[rip + offset_to_previous_saved_ID], random_int
lea rax, [rip + offset_to_relocated__afl_area_ptr]
mov rax, QWORD PTR [rax]
mov rax, QWORD PTR [rax]
mov QWORD PTR[rip + offset_to_saved_shared_memory_address], rax
\end{lstlisting}
The \cref{lst:init_instr} perfomrs three main tasks:
\begin{itemize}
    \item Find the address of forkserver function dereferencing the resolved symbol, store it into \texttt{rax} and the perfomr an indirect call.
    \item Write a random integer (between 0 and 65535\footnote{
        Since the AFL shared map size is 65536 each Basci Block ID must be between 0 and 65535
    }) as previous \gls{BB} ID
    \item Dereference the address of the \gls{SHM} resolved by the dynamic linker and write it into \textit{.mydata} section in 
        order to avoiding dereference it every single time. 
\end{itemize}

\begin{center}
\textbf{Basic Block Instrumntation}    
\end{center}
The instrumentation added for each \gls{BB} could be seen in \cref{lst:bb_instrumentation}.

\begin{lstlisting}[caption={Basic Block Instrumentation},label=lst:bb_instrumentation]
mov QWORD PTR[rip + offset_to_saved_rax], rax
movzx rax, WORD PTR[rip + offset_to_previous_saved_ID]        
xor ax, random_current_basic_block_ID
add rax, QWORD PTR[rip + offset_to_saved_shared_memory_address]
inc BYTE PTR[rax]
mov WORD PTR[rip + offset_to_precedent_bb_saved], right_shifted_current_ID
mov rax, QWORD PTR[rip + offset_to_saved_rax]
\end{lstlisting}

In \cref{lst:bb_instrumentation} all the consideration previously made for the offset computation are still valid. 
The random\_current\_basic\_block\_ID and right\_shifted\_current\_ID are computed by the instrumentator and hardcoded into each \gls{BB}.

For the \gls{BB} instrumentation the algorithm must also takes care of address shifting during instrumentation.
Let's consider the example of \cref{fig:pow_example} and the \gls{BB} address found in that section.
These address was $[\, 0x11a7 , 0x11b5, 0x11bd , 0x11db , 0x11f1 , 0x1253 \,]$
The instrumentation seen in \cref{lst:bb_instrumentation} occupy in total 44 bytes.
Supposing to have instrumented the location 0x11bd, then all the followinig locations must be shifted of 44 bytes.
And this computation must be done for each \gls{IP} found.

But this approach is quite impractical.
GRAPE, in order to overcome the issue of changing addresses, uses \textit{labels}.
It associates, to each \gls{BB}, an unique label. The \gls{BB} could be retrieved by this uniq label univocally.
If the \gls{BB} cahnge starting address, it will not change the label.
So, during the \gls{IP}s search, instead of keep saved the addresses of basic block, the tool keep saved the labels 
of the \gls{BB} to be instrumented, and only after retrieve the addresses of these \gls{BB}.
The final approach can be seen 

\begin{algorithm}[ht]
    \caption{Instrument Basic Block algorithm}
    \label{alg:bb_instr}
    \begin{algorithmic}[1]
        \REQUIRE $bbs=[\, bb\_label_1 \dots bb\_label_n] , mydata\_addr , mygot\_addr$ 
        \ENSURE all the Basic Block are instrumented correctly

        \FOR{$bb\_label \in bbs$}
            \STATE $bb \gets binary.\textsc{getBasciBlockFromLabel}(bb\_label)$ \COMMENT{\small \textcolor{green!60!black}{``binary'' if the main object used by GRAPE and is a virtual representation of the ELF it is rewriting}}
            \STATE $addr \gets bb.initial\_address$
            \STATE $instructions \gets \textsc{computeBasicBlockInstrumentation}(addr,mydata\_addr,mygot\_addr)$ 

            \FOR{$i=\textsc{len}(instructions)$ \TO $1$} 
                    \STATE $instr \gets \textsc{asm}(instructions[\,i\,])$ 
                    \STATE $\textsc{AddInstructionChange}(addr,instr).\textsc{modify}(binary.cfg)$
                    \IF{$\text{.mydata section shifted after .text enlargement}$}
                        \STATE $mydata\_addr \gets mydata.virtual\_address$
                        \STATE $mygot\_addr \gets mygot.virtual\_address$
                        \STATE $instructions \gets \textsc{computeBasicBlockInstrumentation}(addr,mydata\_addr,mygot\_addr)$ 
                        \STATE $\textsc{RemoveInstructionChange}(addr).\textsc{modify}(binary.cfg)$ 
                        \STATE $\textsc{AddInstructionChange}(addr,instr).\textsc{modify}(binary.cfg)$
                    \ENDIF
            \ENDFOR
        \ENDFOR
    \end{algorithmic}
\end{algorithm}

After the instrumentation of each \gls{BB} according to \cref{alg:bb_instr}, the tool will perform 
\textsc{RIPDataCoherenceAnalysis} and \textsc{SymbolCoherenceAnalysis}, adjusting all the shifted references, 
and overwrite the file that \textbf{lief} had previously write.

In the end the tool must check that the \texttt{e\_phoff} field of the \gls{ELF} header, coincide with 
the actual offset of the \textbf{Program Header Table}, and so with the \texttt{p\_offset} field of the 
\textbf{PT\_PHDR} segment.\\
In some binaries, adding new segments (for \textit{.mydata} and \textit{.my\_got}) and adding instructions to other segments, 
could cause the \textbf{Program Header Table} to be shifted and the \texttt{e\_phoff} field to not be modified correctly.

The tool, first, read with \textbf{lief} library the header of the \gls{ELF}, obtaining the \texttt{e\_phoff} and then 
cycle through all the binary segment until it finds the \textbf{PT\_PHDR} segment. When it find the segment, check 
if the two field coincide.
If not, rewrite the byte that represent \texttt{e\_phoff} manually (is the 0x20-th byte of the header), adjusting 
it to the correct value.
The ``manual'' rewriting is perform without an actual library supposing a little-endian addressing.
This is done in order to avoid libraries (such as \textbf{lief}) to adjust other field.
